%%=====================================================================
%%  Surgical fragment: replacement contents for Appendix~\ref{app:lean}
%%  (Machine-verifiable statements).  Insertion-only: no preamble.
%%  Every listing below compiles under Lean 4.15.0 / Mathlib v4.15.0
%%  with zero `sorry` and no custom axioms
%%  (footprint: propext, Classical.choice, Quot.sound).
%%=====================================================================

\noindent The following theorem statements have been discharged in Lean~4
(version 4.15.0) against Mathlib~v4.15.0; each compiles with no
\texttt{sorry} and depends only on the standard foundational axioms
\texttt{propext}, \texttt{Classical.choice}, and \texttt{Quot.sound}.


\medskip \noindent\textbf{A.1\quad Rotation--scaling equivalence (Lemma~\ref{lem:rot}(i)).}
\begin{verbatim}
theorem log_spiral_rotation_scaling
    (a0 b alpha theta : ℝ) :
    a0 * Real.exp (b * (theta - alpha))
      = Real.exp (-(b * alpha)) * (a0 * Real.exp (b * theta)) := by
  rw [Real.exp_neg, mul_sub, Real.exp_sub]
  field_simp
\end{verbatim}

\medskip \noindent\textbf{A.2\quad No common invariance (Lemma~\ref{lem:rot}(ii)).}
\begin{verbatim}
theorem counterwound_no_common_invariance
    (b alpha : ℝ) (hb : b ≠ 0)
    (h : Real.exp (b * alpha) = Real.exp (-(b * alpha))) :
    alpha = 0 := by
  have h2 : b * alpha = -(b * alpha) := Real.exp_injective h
  have h3 : b * alpha = 0 := by linarith
  rcases mul_eq_zero.mp h3 with hb0 | ha0
  · exact absurd hb0 hb
  · exact ha0
\end{verbatim}

\medskip \noindent\textbf{A.3\quad No exchange (Lemma~\ref{lem:rot}(iii)).}
\begin{verbatim}
theorem no_rotation_exchanges_shells
    (a0 b alpha : ℝ) (ha : a0 > 0) (hb : b ≠ 0) :
    ¬ (∀ theta : ℝ,
        Real.exp (-(b * alpha)) * (a0 * Real.exp (b * theta))
          = a0 * Real.exp (-(b * theta))) := by
  intro h
  have h0 := h 0
  have h1 := h 1
  simp only [mul_zero, neg_zero, Real.exp_zero, mul_one] at h0
  have hEbα : Real.exp (-(b * alpha)) = 1 := by
    have hcancel : Real.exp (-(b * alpha)) * a0 = 1 * a0 := by rw [one_mul]; exact h0
    exact mul_right_cancel₀ (ne_of_gt ha) hcancel
  simp only [mul_one, hEbα, one_mul] at h1
  -- h1 : a0 * exp b = a0 * exp (-b)
  have hexp : Real.exp b = Real.exp (-b) :=
    mul_left_cancel₀ (ne_of_gt ha) h1
  have hb0 : b = -b := Real.exp_injective hexp
  have hz : b = 0 := by linarith
  exact hb hz
\end{verbatim}

\medskip \noindent\textbf{A.4\quad Unique stable global minimum of the potential (Lemma~\ref{lem:V}(i)--(iii)).}
\begin{verbatim}
theorem potential_unique_global_min
    (Vp lam bet : ℝ) (hV : Vp > 0) (hl : lam < 0) (hb : bet > 0) :
    ∃! xs : ℝ, xs > 0 ∧
      (∀ x, x > 0 → Vp*(lam*x^4 + bet*x^5) ≥ Vp*(lam*xs^4 + bet*xs^5)) ∧
      xs = (4*(-lam))/(5*bet) := by
  set xs : ℝ := (4*(-lam))/(5*bet) with hxs
  have hbne : bet ≠ 0 := ne_of_gt hb
  have hxs_pos : xs > 0 := by
    rw [hxs]
    apply div_pos
    · linarith
    · linarith
  refine ⟨xs, ⟨hxs_pos, ?_, rfl⟩, ?_⟩
  · -- global minimum: W(x) - W(xs) ≥ 0 for all x > 0
    intro x hx
    -- Reduce to showing lam*x^4+bet*x^5 ≥ lam*xs^4+bet*xs^5, then multiply by Vp>0.
    have hkey : lam*x^4 + bet*x^5 ≥ lam*xs^4 + bet*xs^5 := by
      -- Let g(x) = lam*x^4 + bet*x^5. g'(x) = 4*lam*x^3 + 5*bet*x^4 = x^3(4*lam+5*bet*x).
      -- On (0,∞), g decreasing for x<xs, increasing for x>xs, so xs is the global min.
      -- Prove directly: g(x) - g(xs) ≥ 0 via a factorisation argument.
      -- 5*bet*xs = -4*lam, i.e. 4*lam + 5*bet*xs = 0.
      have hcrit : 4*lam + 5*bet*xs = 0 := by
        rw [hxs]; field_simp; ring
      have hlam : lam = -(5/4)*bet*xs := by linarith
      -- bracket is a sum of nonneg terms for x, xs > 0
      have hbracket : 0 ≤ 4*x^3 + 3*xs*x^2 + 2*xs^2*x + xs^3 := by
        have h1 : 0 ≤ x^3 := by positivity
        have h2 : 0 ≤ xs*x^2 := by positivity
        have h3 : 0 ≤ xs^2*x := by positivity
        have h4 : 0 ≤ xs^3 := by positivity
        linarith
      have hfac :
          lam*x^4 + bet*x^5 - (lam*xs^4 + bet*xs^5)
            = bet * ((1/4)*(x-xs)^2 * (4*x^3 + 3*xs*x^2 + 2*xs^2*x + xs^3)) := by
        rw [hlam]; ring
      have hprod : 0 ≤ bet * ((1/4)*(x-xs)^2 * (4*x^3 + 3*xs*x^2 + 2*xs^2*x + xs^3)) := by
        apply mul_nonneg hb.le
        apply mul_nonneg
        · positivity
        · exact hbracket
      linarith [hfac ▸ hprod]
    have := mul_le_mul_of_nonneg_left hkey (le_of_lt hV)
    linarith
  · -- uniqueness
    rintro y ⟨hy_pos, hy_min, hy_eq⟩
    exact hy_eq
\end{verbatim}

\medskip \noindent\textbf{A.5\quad Reservoir non-negativity and upward bound (Lemma~\ref{lem:V}(iv)--(v)).}
\begin{verbatim}
theorem reservoir_nonneg_and_bounded
    (Vp lam bet : ℝ) (hV : Vp > 0) (hl : lam < 0) (hb : bet > 0) :
    ∀ x, 0 < x → x ≤ (4*(-lam))/(5*bet) →
      0 ≤ (Vp*(lam*x^4 + bet*x^5)
            - Vp*(lam*((4*(-lam))/(5*bet))^4
                  + bet*((4*(-lam))/(5*bet))^5)) := by
  intro x hx hxle
  set xs : ℝ := (4*(-lam))/(5*bet) with hxs
  have hbne : bet ≠ 0 := ne_of_gt hb
  have hxs_pos : xs > 0 := by
    rw [hxs]; exact div_pos (by linarith) (by linarith)
  have hcrit : 4*lam + 5*bet*xs = 0 := by rw [hxs]; field_simp; ring
  have hlam : lam = -(5/4)*bet*xs := by linarith
  have hkey : lam*x^4 + bet*x^5 ≥ lam*xs^4 + bet*xs^5 := by
    have hbracket : 0 ≤ 4*x^3 + 3*xs*x^2 + 2*xs^2*x + xs^3 := by
      have h1 : 0 ≤ x^3 := by positivity
      have h2 : 0 ≤ xs*x^2 := by positivity
      have h3 : 0 ≤ xs^2*x := by positivity
      have h4 : 0 ≤ xs^3 := by positivity
      linarith
    have hfac :
        lam*x^4 + bet*x^5 - (lam*xs^4 + bet*xs^5)
          = bet * ((1/4)*(x-xs)^2 * (4*x^3 + 3*xs*x^2 + 2*xs^2*x + xs^3)) := by
      rw [hlam]; ring
    have hprod : 0 ≤ bet * ((1/4)*(x-xs)^2 * (4*x^3 + 3*xs*x^2 + 2*xs^2*x + xs^3)) := by
      apply mul_nonneg hb.le
      apply mul_nonneg
      · positivity
      · exact hbracket
    linarith [hfac ▸ hprod]
  nlinarith [mul_le_mul_of_nonneg_left hkey (le_of_lt hV)]
\end{verbatim}

\medskip \noindent\textbf{A.6\quad Quarter-turn golden scaling (Corollary~\ref{cor:phi}).}
\begin{verbatim}
theorem quarter_turn_is_phi
    (phi : ℝ) (hphi : phi = (1 + Real.sqrt 5)/2) :
    Real.exp ((2 * Real.log phi / Real.pi) * (Real.pi/2)) = phi := by
  have hpi : Real.pi ≠ 0 := Real.pi_ne_zero
  have hphi_pos : phi > 0 := by
    rw [hphi]
    have h5 : Real.sqrt 5 ≥ 0 := Real.sqrt_nonneg 5
    linarith
  have harg : (2 * Real.log phi / Real.pi) * (Real.pi/2) = Real.log phi := by
    field_simp
  rw [harg, Real.exp_log hphi_pos]
\end{verbatim}

\medskip \noindent\textbf{A.7\quad Ratio invariance (Lemma~\ref{lem:gkf}(ii)).}
\begin{verbatim}
theorem coupling_ratio_warp_independent
    (C Vw Ii Ij : ℝ) (hj : C * Vw * Ij ≠ 0) :
    (C * Vw * Ii) / (C * Vw * Ij) = Ii / Ij := by
  rw [div_eq_div_iff hj (by
        intro hIj
        apply hj
        rw [mul_eq_zero]; right; exact hIj)]
  ring
\end{verbatim}

\medskip \noindent\textbf{A.8\quad Holonomy-neutrality of equal-transport bilinears (Lemma~\ref{lem:neutral}).}
\begin{verbatim}
theorem holonomy_neutral_bilinear
    {n : Type*} [Fintype n] [DecidableEq n]
    (z : ℂ) (hz : Complex.abs z = 1)
    (M : Matrix n n ℂ) (psi : n → ℂ) :
    dotProduct (star (z • psi)) (M.mulVec (z • psi))
      = dotProduct (star psi) (M.mulVec psi) := by
  have hconj : (starRingEnd ℂ) z * z = 1 := by
    have h1 : (starRingEnd ℂ) z * z = ((Complex.normSq z : ℝ) : ℂ) := by
      rw [Complex.normSq_eq_conj_mul_self]
    rw [h1]
    have : Complex.normSq z = 1 := by
      rw [Complex.normSq_eq_abs, hz]; norm_num
    rw [this]; norm_num
  rw [Matrix.mulVec_smul, star_smul, smul_dotProduct, dotProduct_smul, smul_smul]
  simp only [RCLike.star_def]
  rw [hconj, one_smul]
\end{verbatim}

\medskip \noindent\textbf{A.9\quad One-parameter equilibrium-response null (Lemma~\ref{lem:adiabatic}, elementary shadow).}
\begin{verbatim}
theorem cycle_work_null
    (G : ℝ → ℝ) (phi : ℝ → ℝ) (T : ℝ)
    (hG : ContDiff ℝ 1 G)
    (hphi : ContDiff ℝ 1 phi)
    (hper : phi T = phi 0) :
    (∫ t in (0:ℝ)..T, deriv G (phi t) * deriv phi t) = 0 := by
  have hGdiff : Differentiable ℝ G := hG.differentiable (by norm_num)
  have hphidiff : Differentiable ℝ phi := hphi.differentiable (by norm_num)
  -- integrand = deriv (G ∘ phi) by the chain rule
  have hchain : ∀ t, deriv (G ∘ phi) t = deriv G (phi t) * deriv phi t := by
    intro t
    exact deriv_comp t (hGdiff (phi t)) (hphidiff t)
  have hrw : (∫ t in (0:ℝ)..T, deriv G (phi t) * deriv phi t)
        = ∫ t in (0:ℝ)..T, deriv (G ∘ phi) t := by
    apply intervalIntegral.integral_congr
    intro t _; exact (hchain t).symm
  rw [hrw]
  -- G ∘ phi is C¹
  have hGphi : ContDiff ℝ 1 (G ∘ phi) := hG.comp hphi
  have hderiv : ∀ t ∈ Set.uIcc (0:ℝ) T,
      DifferentiableAt ℝ (G ∘ phi) t := by
    intro t _
    exact hGphi.differentiable (by norm_num) t
  have hcont : Continuous (deriv (G ∘ phi)) := by
    have h2 : (deriv (G ∘ phi)) = fun t => deriv G (phi t) * deriv phi t := by
      funext t; exact hchain t
    rw [h2]
    have hcG : Continuous (deriv G) := hG.continuous_deriv (by norm_num)
    have hcphi : Continuous (deriv phi) := hphi.continuous_deriv (by norm_num)
    exact (hcG.comp hphidiff.continuous).mul hcphi
  have hint : IntervalIntegrable (deriv (G ∘ phi)) MeasureTheory.volume 0 T :=
    hcont.intervalIntegrable 0 T
  rw [intervalIntegral.integral_deriv_eq_sub hderiv hint]
  simp only [Function.comp_apply]
  rw [hper]; ring
\end{verbatim}

\medskip \noindent\textbf{A.10\quad Flux-density bound (Proposition~\ref{prop:ceiling}, eq.~\eqref{eq:T10}, core inequality).}
\begin{verbatim}
theorem flux_density_bound (a b : ℝ) :
    a * b ≤ (a^2 + b^2) / 2 := by
  nlinarith [sq_nonneg (a - b)]
\end{verbatim}

\medskip \noindent\textbf{A.11\quad Device-variable reservoir identity and non-negativity (Lemma~\ref{lem:resdev}).}
\begin{verbatim}
theorem reservoir_device_factorisation (u : ℝ) :
    1 - 5*u^4 + 4*u^5
      = (u - 1)^2 * (4*u^3 + 3*u^2 + 2*u + 1) := by
  ring

theorem reservoir_device_nonneg (u : ℝ) (hu : 0 < u) :
    0 ≤ 1 - 5*u^4 + 4*u^5 := by
  nlinarith [sq_nonneg (u - 1), pow_pos hu 3, pow_pos hu 2, hu.le]
\end{verbatim}

\medskip \noindent\textbf{A.12\quad Static-$\NOP$ factorisation (Lemma~\ref{lem:factor}, one-parameter shadow).}
\begin{verbatim}
theorem static_order_parameter_factorisation
    (w : ℝ) (B : ℝ → ℝ) (t : ℝ)
    (hB : DifferentiableAt ℝ B t) :
    deriv (fun s => w * B s) t = w * deriv B t := by
  simp [deriv_const_mul _ hB]
\end{verbatim}

\medskip \noindent\textbf{A.13\quad Coupling-multiplier lower bound (Lemma~\ref{lem:gkexp}).}
\begin{verbatim}
theorem coupling_multiplier_lower_bound (c d : ℝ) (hc : 0 ≤ c) :
    1 + c*d ≤ Real.exp (c*d) := by
  have := Real.add_one_le_exp (c*d); linarith
\end{verbatim}

\medskip \noindent\textbf{B14\quad $\QCD$--electromagnetic channel derivative (Corollary~\ref{cor:channel}).}
\begin{verbatim}
theorem em_channel_derivative
    (g2inv : ℝ → ℝ) (F Phi : ℝ)
    (hg : DifferentiableAt ℝ g2inv Phi) :
    deriv (fun p => (1/4) * g2inv p * F) Phi
      = (1/4) * deriv g2inv Phi * F := by
  have h1 : (fun p => (1/4) * g2inv p * F) = (fun p => ((1/4) * F) * g2inv p) := by
    funext p; ring
  rw [h1, deriv_const_mul _ hg]
  ring
\end{verbatim}

\medskip \noindent\textbf{B16\quad Protection of the electroweak mixing (Corollary~\ref{cor:weinberg}).}
\begin{verbatim}
theorem weinberg_angle_dilaton_independent
    (r : ℝ) :
    deriv (fun _ : ℝ => (1 + r)⁻¹) = fun _ => 0 := by
  funext Phi
  simp
\end{verbatim}
